% ===================== preambolo =====================
\documentclass[11pt,a4paper]{article}

\usepackage[T1]{fontenc}
\usepackage[utf8]{inputenc}
\usepackage[italian]{babel}
\usepackage[a4paper,margin=25mm]{geometry}
\usepackage{amsmath,amssymb}
\usepackage{graphicx}
\usepackage{booktabs}
\usepackage{multirow}
\usepackage{xcolor}
\usepackage{caption}
\usepackage{enumitem}
\usepackage{siunitx}
\sisetup{per-mode=symbol,range-phrase={--},range-units=single}
\usepackage[hidelinks]{hyperref}

\captionsetup{font=small,labelfont=bf}
\setlist{itemsep=2pt,parsep=0pt,topsep=4pt}

% --- segnaposto per figure: compila senza immagini e si vede cosa manca ---
% Uso:  \figph{55mm}{fig:etichetta}{Didascalia}{che cosa deve mostrare}
\newcommand{\figph}[4]{%
  \begin{figure}[htbp]
    \centering
    \fcolorbox{gray!60}{gray!5}{%
      \begin{minipage}[c][#1][c]{0.92\linewidth}
        \centering\ttfamily\footnotesize\color{gray!80!black}
        [SEGNAPOSTO FIGURA]\\[2pt]
        #4
      \end{minipage}}
    \caption{#3}
    \label{#2}
  \end{figure}}

% abbreviazioni dei controllori
\newcommand{\Co}{\textsf{C1}}
\newcommand{\Ct}{\textsf{C2}}
\newcommand{\Cth}{\textsf{C3}}

% =====================================================================
\title{\textbf{Controllo con interazione col terreno di un robot esapode:\\
confronto misurato di tre controllori\\ su un impianto di simulazione condiviso}}

% TODO: inserire i nomi veri e le matricole
\author{[Studente 1] \and [Studente 2]}
\date{Field and Service Robotics --- A.A. 2025/2026\\
Universit\`a degli Studi di Napoli Federico II\\[4pt]
\today}

\begin{document}
\maketitle

\begin{abstract}
\noindent
Confrontiamo tre controllori per il PhantomX AX Metal Hexapod Mark~II su un
unico impianto Simscape Multibody ad alta fedelt\`a: un controllore cinematico
a tripode ad anello aperto (\Co), la retroazione di ricerca del terreno basata
sulla coppia di Arrigoni et al.~\cite{arrigoni2024} (\Ct) e un anello
d'assetto aggiunto sopra di essa (\Cth). Ogni controllore aggiunge al
precedente esattamente un meccanismo, cos\`i che ogni confronto fra controllori
adiacenti isoli una sola causa.

Il risultato principale del progetto non era nel piano iniziale. Circa met\`a
del lavoro \`e consistita nel \emph{trovare e correggere due difetti del nostro
stesso banco di prova}, entrambi capaci di invertire le conclusioni e nessuno
dei due dichiarato: il robot camminava, l'animazione era plausibile e le
metriche uscivano. Li riportiamo per intero, insieme alla disciplina di misura
adottata in seguito --- un rumore di fondo misurato, una soglia di leggibilit\`a
fissata prima di guardare i dati e ritrattazioni scritte --- che consideriamo il
contributo pi\`u trasferibile di questo lavoro.

Con quella disciplina applicata, il risultato misurato \`e che la ricerca del
terreno \emph{costa} assetto dove non c'\`e nulla da cercare e lo \emph{ripaga}
sulle discontinuit\`a, mentre l'anello d'assetto restituisce quell'assetto e lo
paga in energia. Nessuno dei due meccanismi cambia la risposta a un impulso
laterale e nessuno dei tre completa il percorso a sette ostacoli.
\end{abstract}

\tableofcontents
\newpage

% =====================================================================
\section{Introduzione}
\label{sec:intro}

\subsection{La piattaforma e il lavoro di riferimento}

Il PhantomX AX Metal Hexapod Mark~II \`e un robot esapode economico e a
hardware aperto, con diciotto servomotori Dynamixel AX-12A, tre per zampa
(coxa, femore, tibia). Nella configurazione standard \`e privo di sensori di
forza ai piedi e di IMU: l'unica retroazione disponibile dall'hardware \`e la
posizione e la stima di carico riportate da ciascun servo.

Arrigoni et al.~\cite{arrigoni2024} sfruttano esattamente questo vincolo. Il
loro contributo \`e un'architettura di retroazione che rileva il contatto
piede--terreno \emph{dalla sola coppia motrice}, confrontando la coppia
misurata ai giunti con quella prevista da un modello di dinamica inversa del
robot, e usa la rilevazione per adattare l'altezza comandata del piede a un
terreno che il generatore di traiettoria non conosce. L'attrattiva dell'idea
\`e che compra adattamento al terreno senza aggiungere un solo sensore.

\subsection{Perch\'e abbiamo scelto questo progetto}

Tre ragioni hanno reso questo un progetto finale interessante, e non un
esercizio di reimplementazione.

\begin{description}[style=nextline,leftmargin=0pt]
\item[L'affermazione \`e falsificabile, e l'articolo non ne quantifica il prezzo.]
Il lavoro di riferimento dimostra che la ricerca del terreno migliora il
comportamento su terreno non modellato. Non riporta quanto costi il meccanismo
dove non c'\`e terreno da scoprire. Un confronto controllato su terreno
identico, con un criterio di leggibilit\`a dichiarato, pu\`o rispondere --- e la
risposta si \`e rivelata essere un costo reale e sistematico
(Sezione~\ref{sec:results}).

\item[Il rilevatore di contatto \`e un osservatore basato su modello travestito.]
La flag che muove l'intera architettura \`e
$\mathrm{OR}_j\left(|\tau_j^{\mathrm{mis}} - \tau_j^{\mathrm{att}}| >
\sigma_j\right)$, dove $\tau^{\mathrm{att}}$ proviene da un modello di dinamica
inversa a struttura di corpi rigidi. La sua qualit\`a dipende quindi
dall'\emph{accuratezza di un modello dinamico}, il che si collega direttamente
al materiale di modellistica e stima dello stato del corso. Abbiamo scoperto a
nostre spese quanto direttamente (Sezione~\ref{sec:estimator}).

\item[Un esapode su un impianto condiviso rende pulita l'ablazione.]
Poich\'e tutti e tre i controllori pilotano lo stesso modello multibody, lo
stesso terreno, la stessa legge di contatto e lo stesso generatore di andatura,
ogni confronto fra adiacenti differisce per un solo meccanismo. Questo \`e
raramente possibile quando si confrontano fra loro controllori pubblicati.
\end{description}

\subsection{Che cosa sostiene questo report}

\begin{enumerate}
\item Due difetti del banco di simulazione, trovati, misurati, corretti e
      documentati, ciascuno dei quali aveva gi\`a prodotto numeri dall'aria
      pubblicabile ma sbagliati (Sezione~\ref{sec:defects}).
\item Un protocollo di misura che rende una differenza fra due controllori
      \emph{dichiarabile o non dichiarabile}, invece che semplicemente
      osservata (Sezione~\ref{sec:method}).
\item Un confronto a tre su sette task con quel protocollo applicato
      (Sezione~\ref{sec:results}).
\item Una verifica di fattibilit\`a delle coppie comandate rispetto al datasheet
      del servo reale (Sezione~\ref{sec:feasibility}).
\end{enumerate}

% =====================================================================
\section{Impostazione sperimentale}
\label{sec:setup}

\subsection{L'impianto}

Il robot \`e importato in Simscape Multibody dall'URDF pubblico del pacchetto di
descrizione del PhantomX: 25 corpi rigidi (tronco pi\`u 24 membri delle zampe),
18 giunti rotoidali e sei elementi di contatto sferici ai piedi. I giunti sono
comandati in \emph{posizione}; la coppia di giunto che ne risulta \`e quindi una
reazione, non un ingresso --- un punto che conta per la
Sezione~\ref{sec:feasibility}. L'interazione piede--terreno usa una legge di
contatto a penalit\`a con attrito di Coulomb, tarata in modo che la penetrazione
statica ($4.3$~mm in appoggio a tripode) stia ben sopra la regione di
transizione ($2.0$~mm), dove altrimenti la legge di forza sarebbe mal definita.

\figph{60mm}{fig:plant}{L'impianto Simscape Multibody: corpo, sei zampe a tre gradi di libert\`a, elementi di contatto sferici e i terreni usati dalla batteria di task.}{screenshot del modello Simscape / schema a blocchi dell'impianto}

\subsection{I tre controllori}

\begin{table}[htbp]
\centering
\caption{I tre controllori. Ciascuno aggiunge un meccanismo al precedente.}
\label{tab:controllers}
\small
\begin{tabular}{@{}llp{85mm}@{}}
\toprule
 & \textbf{Impianto} & \textbf{Meccanismo} \\
\midrule
\Co  & \texttt{sim\_zero} & Tripode cinematico ad anello aperto. Traiettorie
  dei piedi generate dai parametri d'andatura, cinematica inversa per zampa,
  giunti comandati in posizione. Nessuna retroazione di alcun tipo. \\[3pt]
\Ct  & \texttt{sim\_zero} & \Co{} pi\`u la \emph{ricerca del terreno}: finch\'e
  una zampa \`e in appoggio e la flag di contatto \`e bassa, l'altezza comandata
  del piede continua a scendere fino a che il residuo di coppia supera la
  soglia. \\[3pt]
\Cth & \texttt{sim\_attitude} & \Ct{} pi\`u un \emph{anello d'assetto}: due
  regolatori PI discreti su beccheggio e rollio del corpo, le cui uscite sono
  distribuite sui sei piedi come scostamento additivo di altezza. \\
\bottomrule
\end{tabular}
\end{table}

L'anello d'assetto agisce \emph{in serie, dopo} la ricerca del terreno: legge
l'altezza del piede gi\`a corretta e vi somma un termine. Lo abbiamo verificato
leggendo la connettivit\`a dei blocchi dentro il file del modello, non dalla
documentazione:

\begin{center}
\texttt{ricerca\_terreno} $\rightarrow z_i \rightarrow \Sigma(+\Delta z_i)
\rightarrow$ \texttt{Saturation} $\rightarrow$ \texttt{Rate Limiter}
$\rightarrow$ zampa
\end{center}

con
\begin{equation}
\Delta z = \mathbf{x}\,u_\theta + \mathbf{y}\,u_\phi ,
\label{eq:dz}
\end{equation}
dove $u_\theta$ e $u_\phi$ sono le uscite dei regolatori di beccheggio e rollio
e $\mathbf{x},\mathbf{y}\in\mathbb{R}^6$ sono i bracci di leva dei sei piedi.
Il beccheggio usa $P=10$, $I=1$; il rollio $P=5$, $I=0.5$; entrambi sono
discreti a \SI{1}{\milli\second} con riferimento nullo. Poich\'e \Cth{} contiene
\Ct{}, il confronto $\Ct\rightarrow\Cth$ isola l'anello d'assetto e
$\Co\rightarrow\Ct$ isola la ricerca del terreno.

\subsection{Parametri dell'andatura e loro provenienza}
\label{sec:gait}

Il lavoro di riferimento delega la generazione dell'andatura a NUKE e dichiara
esplicitamente che non la discuter\`a. I parametri d'andatura andavano quindi
scelti. Non sono arbitrari: tre su cinque sono fissati da un vincolo misurato,
e il grado di libert\`a residuo \`e uno solo, perch\'e
\begin{equation}
T = \frac{S}{\beta_{\mathrm{app}}\, v_{\mathrm{nom}}} .
\label{eq:gait}
\end{equation}

\begin{table}[htbp]
\centering
\caption{Parametri dell'andatura e loro provenienza.}
\label{tab:gait}
\small
\begin{tabular}{@{}llp{78mm}@{}}
\toprule
\textbf{Parametro} & \textbf{Valore} & \textbf{Origine} \\
\midrule
vettore delle fasi & --- &
  Definizione del tripode alternato: gruppi
  $\{\mathrm{AS,CD,PS}\}$ e $\{\mathrm{AD,CS,PD}\}$. Non \`e un parametro
  libero. \\
$\beta_{\mathrm{app}}$ & $0.50$ &
  L'unico fattore di servizio che d\`a appoggio continuo su \emph{esattamente}
  tre zampe. Sotto, esistono istanti con meno di tre piedi a terra (fuori dalla
  stabilit\`a statica); sopra, compaiono finestre a sei piedi in cui 18 vincoli
  di posizione agiscono su 6 gradi di libert\`a. \\
$S$ & \SI{0.06}{\meter} &
  Limitato dall'estensione della zampa. A met\`a passo la zampa usa l'$80.9\%$
  di $\ell_f+\ell_t=\SI{0.219}{\meter}$; l'assert interno scatta all'$85\%$,
  cio\`e a $S=\SI{0.095}{\meter}$. Il limite dell'$85\%$ \`e a sua volta
  misurato: all'$89\%$ il robot camminava all'indietro. \\
$H$ & \SI{0.05}{\meter} &
  $1.5\times$ l'altezza misurata dell'ostacolo (\SI{32.8}{\milli\meter}, letta
  dalla quota del piede in appoggio nella run a ostacolo singolo, non dalla
  mesh). \\
$T$ & \SI{1.00}{\second} &
  Variabile dipendente, tramite la~\eqref{eq:gait}. D\`a
  $v_{\mathrm{nom}}=\SI{0.12}{\meter\per\second}$, verificata all'$83\%$
  dell'inviluppo di stabilit\`a misurato (Sezione~\ref{sec:envelope}). \\
\bottomrule
\end{tabular}
\end{table}

Il numero di Froude che ne risulta \`e
$v^2/(g h) = 0.12^2/(9.81\times 0.154) = 0.0095$: le forze inerziali sono
l'uno percento della gravit\`a. La marcia \`e nettamente quasi-statica, il che
\`e ci\`o che legittima l'ipotesi di stabilit\`a \emph{statica} su cui \`e
costruita l'andatura a tripode.

\subsection{Batteria di task}

\begin{table}[htbp]
\centering
\caption{I sette task.}
\label{tab:tasks}
\small
\begin{tabular}{@{}llp{80mm}@{}}
\toprule
\textbf{ID} & \textbf{Nome} & \textbf{Che cosa sollecita} \\
\midrule
T2  & piano, sweep di velocit\`a & Riferimento, e inviluppo di stabilit\`a in
                          velocit\`a. \\
T3  & percorso in curva & Comando di imbardata \SI{0.1}{\radian\per\second}. \\
T4  & rampa di \SI{8}{\degree} & Pendenza continua: assetto del corpo contro
                          terreno. \\
T4D & dosso \SI{8}{\degree}/\SI{8}{\degree} & Salita e poi discesa, due
                          interfacce angolari. \\
T5  & ostacolo singolo   & Un gradino non modellato, \SI{32.8}{\milli\meter}. \\
T6  & sette ostacoli     & Discontinuit\`a non modellate ripetute, a mezza
                          larghezza perch\'e il lato sinistro e il destro siano
                          sollevati alternativamente. \\
T7  & impulso laterale   & Gradino di quantit\`a di moto dall'$1\%$ al
                          $1600\%$ del nominale. \\
\bottomrule
\end{tabular}
\end{table}

\figph{65mm}{fig:tasks}{Le configurazioni di terreno della batteria di task.}{montaggio dei terreni: piano, rampa, dosso, ostacolo singolo, percorso a sette ostacoli}

% =====================================================================
\section{Metodo: quando una differenza \`e un risultato?}
\label{sec:method}

Tre regole, ciascuna adottata dopo averla violata almeno una volta. Le
consideriamo la parte pi\`u trasferibile di questo lavoro, perch\'e sono ci\`o che
ha separato i risultati della Sezione~\ref{sec:results} dai numeri in cui
credevamo prima.

\subsection{Un rumore di fondo misurato, e una soglia di leggibilit\`a}
\label{sec:noise}

Ogni run \`e ripetuta tre volte con l'altezza di appoggio $z_0$ spostata di
$\pm\SI{0.2}{\milli\meter}$ --- una quantit\`a che non ha nulla a che vedere con
il controllore e tutto a che vedere con dove capita che sia il terreno.
L'escursione che ne risulta per ciascuna metrica \`e il \emph{rumore di fondo}
di quella metrica, su quel task, per quel controllore.

\begin{quote}
Una differenza fra due controllori si dichiara solo se \`e almeno
\textbf{tre volte} il peggiore dei due rumori di fondo.
\end{quote}

Non \`e una formalit\`a. Applicarla ha eliminato diverse differenze che
sembravano risultati; applicarla \emph{male} --- con un rumore misurato su un
controllore diverso --- ci ha fatto credere per un giorno che non fosse
dichiarabile proprio nulla. Il rumore \`e misurato su quattro task (piano,
curva, dosso, ostacolo) per tutti e tre i controllori, ciascuno contro s\'e
stesso.

In tutta la Sezione~\ref{sec:results}, una cifra fra parentesi come
$(16.3\times)$ \`e il rapporto fra la differenza osservata e quel rumore di
fondo. Le voci marcate \textsf{n.c.} (non concludente) stanno sotto il
$3\times$; quelle marcate \textsf{n.m.} non hanno rumore misurato su quel task
e sono riportate senza verdetto.

\subsection{Il criterio si scrive prima della run}

Ogni script diagnostico del repository porta in intestazione la domanda che
pone, la soglia che applicher\`a e le conclusioni possibili --- tutto scritto
prima che i dati esistano. Questo conta perch\'e l'alternativa \`e disponibile e
allettante: in questo progetto una soglia di leggibilit\`a calibrata dopo aver
visto i dati avrebbe spostato il limite di stabilit\`a da $1.0\times$ a
$1.2\times$ a seconda del multiplo scelto.

\subsection{Le ritrattazioni restano scritte}

Ogni affermazione ritirata \`e marcata
\texttt{[RITRATTATA \textless data\textgreater]} sul posto, con la motivazione.
Ne abbiamo ritirate almeno cinque nel corso del progetto. In questo lavoro
un'affermazione sbagliata \`e sopravvissuta tre giorni proprio perch\'e era
plausibile, e la disciplina esiste per rendere visibile questo, non per
mettere in ordine.

% =====================================================================
\section{Due difetti del banco di prova}
\label{sec:defects}

Il blocco di lavoro singolo pi\`u grande di questo progetto non \`e stato il
confronto. \`E stato scoprire che il banco con cui stavamo misurando era
sbagliato in due punti, e che il secondo difetto era \emph{causato dalla
correzione del primo}.

\subsection{Le inerzie dell'URDF sono sbagliate di circa un fattore mille}
\label{sec:inertias}

L'impianto prende masse e inerzie da un URDF pubblico. Le masse sono corrette;
i tensori d'inerzia no. Due argomenti indipendenti, nessuno dei quali dipende
dal fidarsi di una particolare fonte per il file:

\begin{description}[style=nextline,leftmargin=10pt]
\item[Raggio d'inerzia.] Da $I = m r^2$, l'inerzia dichiarata del tronco
colloca la massa del tronco a $r = \SI{1.79}{\meter}$ dal proprio centro. Il
tronco \`e largo \SI{25}{\centi\meter}. Ogni membro d\`a
$r = \SI{0.46}{\meter}$; i membri sono lunghi \SIrange{2}{12}{\centi\meter}.
Letti in $\mathrm{kg}\cdot\mathrm{m}^2$, come lo standard URDF richiede, quei
numeri descrivono un oggetto che non pu\`o esistere.
\item[Tutti i 24 membri portano inerzia identica,] allo stesso decimale. Coxa,
femore e tibia differiscono per forma e lunghezza. \`E un segnaposto, non una
misura.
\end{description}

Dividere per $10^3$ riporta tutto nell'intervallo fisicamente plausibile, che
\`e la firma di un export CAD eseguito con le masse in grammi. L'effetto sulle
conclusioni, e non solo sui valori, \`e grande: con le inerzie dell'URDF il
beccheggio di \Co{} su terreno accidentato raggiungeva \SI{45}{\degree}, contro
\SI{15}{\degree} con quelle corrette, e il rapporto di costo fra \Ct{} e \Co{}
passava da $1.35$ a $0.68$ --- \emph{cambiava segno}.

\subsection{Correggere il robot ha rotto il rilevatore di contatto}
\label{sec:estimator}

La flag di contatto di \Ct{} non \`e una soglia sulla coppia misurata. Il
modello calcola
\begin{equation}
\mathrm{contatto} = \bigvee_{j\in\{c,f,t\}}
\left( \left| \tau_j^{\mathrm{mis}} - \tau_j^{\mathrm{att}} \right| > \sigma_j \right),
\label{eq:flag}
\end{equation}
e $\tau^{\mathrm{att}}$ \`e prodotta da un blocco di dinamica inversa che porta
un \texttt{rigidBodyTree} importato \emph{dallo stesso URDF}.

Avevamo corretto le inerzie del solo robot simulato. Da quel momento il robot e
il modello che ne prevedeva la coppia erano due robot diversi, con inerzie
distanti tre ordini di grandezza. Il residuo
$|\tau^{\mathrm{mis}} - \tau^{\mathrm{att}}|$ \`e diventato grande
\emph{sempre}, con o senza contatto, e la flag \`e rimasta alta per oltre il
$95\%$ della fase di volo. Con la flag bloccata a uno, la ricerca del terreno
non entra mai nel ramo che fa crescere l'estensione del piede, e il comando
degenera in una costante.

\begin{quote}
Per tre giorni \Ct{} non ha cercato il terreno. Nulla lo segnalava: il robot
camminava, l'animazione era plausibile, le metriche uscivano. Ogni riga \Ct{}
prodotta in quella finestra descriveva un controllore rotto.
\end{quote}

\paragraph{Come l'abbiamo misurato.} Una singola run di dieci secondi, in cui
il log contiene contemporaneamente la coppia misurata, la coppia attesa e la
\emph{forza normale al piede}. La flag pu\`o allora essere ricalcolata offline
per qualunque soglia e confrontata con la forza, che \`e la verit\`a di
riferimento.

\begin{table}[htbp]
\centering
\caption{Flag di contatto contro forza misurata al piede, a soglia invariata.
La seconda riga \`e la diagnosi: una flag sempre alta sbaglia esattamente tante
volte quanto la nostra, quindi la nostra non portava informazione.}
\label{tab:estimator}
\small
\begin{tabular}{@{}lcc@{}}
\toprule
 & \textbf{stimatore disallineato} & \textbf{allineato} \\
\midrule
errore della flag durante l'appoggio      & $7.0\%$ & $\mathbf{0.6\%}$ \\
errore di una flag \emph{sempre alta}     & $7.0\%$ & $3.2\%$ \\
la flag porta informazione?               & \textbf{no} & s\`i \\
voli senza alcun reset della ricerca      & $20/60$ & $2/60$ \\
\bottomrule
\end{tabular}
\end{table}

\paragraph{La correzione.} Una routine riscrive massa e inerzia di ogni corpo
dell'albero di corpi rigidi dello stimatore con gli stessi valori corretti
applicati all'impianto, forza il sottosistema mascherato a rivalutarsi e
verifica in tre modi successivi --- scrittura, fallback, rilettura --- fermandosi
con errore se uno qualunque fallisce. Un fallimento silenzioso qui
ripristinerebbe esattamente il problema che chiude. \`E chiamata subito dopo la
correzione delle inerzie in ogni script di task, in memoria, senza mai salvare
il modello.

\subsection{Nuova taratura della soglia di contatto}
\label{sec:threshold}

Con lo stimatore allineato la flag funziona, ma alla soglia originale
$[0.04,\,0.1,\,0.1]$ resta comunque alta per oltre il $41.9\%$ del volo, e il
$39.7\%$ di quello \`e la sola coxa: il suo $p_{90}$ in volo ($0.0421$) si
sovrappone al suo $p_{10}$ in appoggio ($0.0469$). Questo ha una spiegazione
fisica --- la coxa ruota attorno all'asse verticale mentre il contatto \`e
verticale, quindi il contatto non le trasmette quasi nessuna coppia, mentre in
volo porta tutta l'accelerazione laterale della zampa.

La terna adottata \`e $[0.385,\,0.136,\,0.105]$, che pone la soglia della coxa
sopra la sua stessa mediana in appoggio, cio\`e la disattiva. \textbf{Questo \`e
un cambio di regola, non una ritaratura:} la flag diventa un OR su due giunti.
\`E stata decisa sul percorso a ostacoli, dove l'esito \`e binario e non soggetto
al rumore di misura, e va dichiarata come nostra scelta: non \`e pi\`u il valore
presente nell'implementazione di riferimento.

\subsection{Un limitatore di pendenza, e perch\'e sta nell'impianto}
\label{sec:slew}

Il modello con assetto portava in origine un limitatore di pendenza di
$\pm\SI{0.3}{\meter\per\second}$ su ciascun comando di altezza del piede, che
l'impianto ad anello aperto non aveva. \`E un'asimmetria: il limitatore vincola
\emph{ogni} comando di altezza, compreso quello prodotto dalla ricerca del
terreno, quindi un confronto \Ct{}--\Cth{} avrebbe confuso l'anello d'assetto
con il suo stesso limitatore.

Abbiamo quindi messo lo stesso limitatore, a
$\pm\SI{0.4}{\meter\per\second}$, in entrambi gli impianti. Il valore non \`e
libero: l'andatura nominale ha un picco di \SI{0.375}{\meter\per\second} di
velocit\`a verticale del piede, quindi a $\pm0.4$ la marcia nominale resta
intatta e il limitatore agisce solo sui comandi rapidi --- ricerca del terreno e
correzioni d'assetto. Abbiamo verificato che le righe \Co{} e il rumore di
fondo di \Co{} sono identici byte per byte con e senza limitatore su cinque dei
sette task e sulla cella nominale del sesto.

Si \`e rivelato pi\`u importante del previsto: il limitatore ha eliminato la
deriva laterale di \Ct{} sul percorso a ostacoli (da $0.296$ a
\SI{0.104}{\meter}) e un'anomala sensibilit\`a al rumore dell'energia di \Cth{}.
Un dettaglio implementativo, non un controllore, con effetti pi\`u grandi di
alcune delle differenze fra controllori.

% =====================================================================
\section{Risultati}
\label{sec:results}

Tutte le righe che seguono vengono da un'unica campagna eseguita dopo le
correzioni della Sezione~\ref{sec:defects}, con il limitatore di pendenza
presente su tutti e tre gli impianti: sette task $\times$ tre controllori, pi\`u
un rumore di fondo di quattro task $\times$ tre run $\times$ tre controllori. I
verdetti in grassetto superano la regola del $3\times$ della
Sezione~\ref{sec:noise}.

\subsection{Piano e curva: la ricerca del terreno costa assetto}

\begin{table}[htbp]
\centering
\caption{Piano (T2) e percorso in curva (T3), velocit\`a nominale. Fra parentesi
il rapporto sul rumore di fondo; \textsf{n.c.} = sotto il $3\times$.}
\label{tab:t2t3}
\small
\begin{tabular}{@{}lrrrll@{}}
\toprule
 & \Co & \Ct & \Cth & $\Co\rightarrow\Ct$ & $\Ct\rightarrow\Cth$ \\
\midrule
\multicolumn{6}{@{}l}{\emph{T2 --- piano, velocit\`a nominale}}\\
beccheggio max [deg]      & 0.514 & 1.95  & 0.356 & \textbf{+279\%} (16.3) & \textbf{$-$82\%} (18.1) \\
rollio max [deg]          & 0.206 & 0.280 & 0.205 & +36\% (n.c.)           & $-$27\% (n.c.) \\
costo di trasporto        & 0.888 & 0.865 & 1.11  & $-$3\% (n.c.)          & \textbf{+28\%} (27.7) \\
energia [J]               & 18.6  & 18.2  & 23.3  & $-$2\% (n.c.)          & \textbf{+28\%} (28.4) \\
coppia RMS [N\,m]         & 0.400 & 0.403 & 0.372 & +1\% (n.c.)            & \textbf{$-$8\%} (8.0) \\
coppia di picco [N\,m]    & 1.87  & 2.02  & 2.23  & +8\% (n.c.)            & +11\% (n.c.) \\
velocit\`a media [m/s]      & 0.135 & 0.136 & 0.135 & \textbf{+1\%} (6.9)    & $-$0\% (n.c.) \\
\midrule
\multicolumn{6}{@{}l}{\emph{T3 --- curva, $+0.1$ rad/s}}\\
beccheggio max [deg]      & 0.579 & 1.85  & 0.392 & \textbf{+219\%} (40.0) & \textbf{$-$79\%} (45.9) \\
rollio max [deg]          & 0.224 & 0.351 & 0.210 & \textbf{+57\%} (4.6)   & \textbf{$-$40\%} (5.1) \\
costo di trasporto        & 0.970 & 0.951 & 1.22  & \textbf{$-$2\%} (3.6)  & \textbf{+28\%} (24.3) \\
energia [J]               & 30.3  & 29.5  & 37.8  & \textbf{$-$3\%} (6.4)  & \textbf{+28\%} (22.5) \\
coppia RMS [N\,m]         & 0.398 & 0.389 & 0.376 & \textbf{$-$2\%} (7.0)  & \textbf{$-$3\%} (7.3) \\
coppia di picco [N\,m]    & 1.91  & 2.11  & 2.43  & +10\% (n.c.)           & +15\% (n.c.) \\
\bottomrule
\end{tabular}
\end{table}

Su un terreno che \`e esattamente quello che il generatore di traiettoria crede
che sia, la ricerca del terreno rende l'assetto del corpo \emph{peggiore} di un
fattore tre o quattro, e non compra nulla: energia e costo di trasporto restano
invariati entro il rumore sul piano e migliorano del due o tre percento in
curva. Il meccanismo continua a lavorare dove non c'\`e nulla da trovare, e il
corpo lo paga.

\figph{70mm}{fig:pitch-flat}{Beccheggio del corpo su due cicli d'andatura in piano, per i tre controllori. L'oscillazione di \Ct{} \`e la ricerca del terreno che agisce su un suolo che gi\`a conosce.}{grafico beccheggio vs tempo, T2, tre curve sovrapposte}

\subsection{Discontinuit\`a: la ricerca del terreno ripaga}

\begin{table}[htbp]
\centering
\caption{Dosso (T4D) e ostacolo singolo (T5).}
\label{tab:t4dt5}
\small
\begin{tabular}{@{}lrrrll@{}}
\toprule
 & \Co & \Ct & \Cth & $\Co\rightarrow\Ct$ & $\Ct\rightarrow\Cth$ \\
\midrule
\multicolumn{6}{@{}l}{\emph{T4D --- dosso, \SI{8}{\degree}/\SI{8}{\degree}}}\\
beccheggio max [deg] & 8.65  & 10.4  & 1.98  & \textbf{+21\%} (19.8)  & \textbf{$-$81\%} (93.8) \\
rollio max [deg]     & 2.12  & 1.49  & 0.611 & \textbf{$-$30\%} (6.4) & \textbf{$-$59\%} (7.6) \\
costo di trasporto   & 1.12  & 1.04  & 1.45  & \textbf{$-$7\%} (5.1)  & \textbf{+40\%} (16.8) \\
energia [J]          & 64.1  & 59.9  & 83.6  & \textbf{$-$7\%} (3.5)  & \textbf{+40\%} (18.8) \\
velocit\`a media [m/s] & 0.130 & 0.133 & 0.131 & \textbf{+2\%} (3.6)    & \textbf{$-$1\%} (3.0) \\
\midrule
\multicolumn{6}{@{}l}{\emph{T5 --- ostacolo singolo, \SI{32.8}{\milli\meter}}}\\
beccheggio max [deg] & 7.56  & 6.41  & 1.73  & \textbf{$-$15\%} (16.4)& \textbf{$-$73\%} (19.8) \\
rollio max [deg]     & 4.49  & 2.31  & 1.13  & $-$49\% (n.c.)         & $-$51\% (n.c.) \\
costo di trasporto   & 1.20  & 1.08  & 1.40  & \textbf{$-$10\%} (5.2) & \textbf{+30\%} (17.1) \\
energia [J]          & 48.1  & 44.3  & 57.5  & \textbf{$-$8\%} (3.6)  & \textbf{+30\%} (17.2) \\
coppia RMS [N\,m]    & 0.408 & 0.403 & 0.380 & $-$1\% (n.c.)          & \textbf{$-$6\%} (9.9) \\
coppia di picco [N\,m] & 4.46 & 8.33 & 4.01  & +87\% (n.c.)           & $-$52\% (n.c.) \\
ostacolo superato    & s\`i   & s\`i   & s\`i   & uguale                 & uguale \\
\bottomrule
\end{tabular}
\end{table}

Qui il quadro si rovescia, ed \`e questo il risultato che riproduce
l'affermazione del lavoro di riferimento: su una discontinuit\`a la ricerca
migliora rollio, energia, costo di trasporto e completamento del task. Resta
una solida eccezione sul dosso: \textbf{il picco di beccheggio \`e peggiore del
$21\%$}, a un rapporto di $19.8$. Su una pendenza continua la ricerca estende i
piedi a valle e il corpo segue il gradiente --- si adatta al terreno invece di
opporvisi.

\figph{70mm}{fig:obstacle}{Assetto del corpo durante il superamento dell'ostacolo (T5). In ombreggiato: la finestra in cui almeno un piede \`e sull'ostacolo, identificata cinematicamente.}{grafico assetto vs tempo su T5, con finestra del passaggio evidenziata}

\subsection{La rampa: a che cosa serve l'anello d'assetto}

\begin{table}[htbp]
\centering
\caption{Rampa di \SI{8}{\degree} (T4). Rumore di fondo non misurato su questo
task; non si emettono verdetti, ma il comportamento \`e qualitativo e
inequivocabile.}
\label{tab:t4}
\small
\begin{tabular}{@{}lrrr@{}}
\toprule
 & \Co & \Ct & \Cth \\
\midrule
assetto del corpo meno angolo della rampa [deg] & 0.106 & 0.259 & $-7.54$ \\
beccheggio picco-picco sulla rampa [deg]        & 0.994 & 1.97  & 0.839 \\
rollio sulla rampa [deg]                        & 0.662 & 0.563 & 0.314 \\
costo di trasporto                              & 1.15  & 1.09  & 1.58 \\
salita riuscita                                 & s\`i   & s\`i   & s\`i \\
\bottomrule
\end{tabular}
\end{table}

\Co{} e \Ct{} allineano il corpo alla rampa, entro qualche decimo di grado.
\Cth{} lo tiene a $7.54^\circ$ dalla rampa, cio\`e entro $0.5^\circ$
dall'\emph{orizzontale}, che \`e esattamente ci\`o per cui \`e costruito un anello
con riferimento d'assetto nullo --- ed \`e anche il comportamento che il lavoro
di riferimento dichiara come proprio obiettivo. Costa il $45\%$ in costo di
trasporto.

\figph{70mm}{fig:ramp}{Beccheggio del corpo durante la salita della rampa di \SI{8}{\degree}. \Co{} e \Ct{} seguono la pendenza; \Cth{} tiene il tronco orizzontale.}{grafico beccheggio vs tempo su T4, tre curve, con la pendenza della rampa come riferimento}

\subsection{Percorso a ostacoli: il confronto non discrimina}

\begin{table}[htbp]
\centering
\caption{Percorso a sette ostacoli (T6). Si riportano solo l'esito binario e la
sua validit\`a.}
\label{tab:t6}
\small
\begin{tabular}{@{}lrrr@{}}
\toprule
 & \Co & \Ct & \Cth \\
\midrule
distanza percorsa [m]        & 2.91  & 3.12  & 2.99 \\
deriva laterale massima [m]  & 0.067 & 0.104 & 0.058 \\
percorso superato, valido    & no    & no    & no \\
\bottomrule
\end{tabular}
\end{table}

``Superato'' vale solo se la deriva laterale massima resta sotto
\SI{0.183}{\meter}; oltre, il robot pu\`o aver mancato per intero un ostacolo
invece di scavalcarlo. La soglia viene dalla geometria degli ostacoli, calcolata
indipendentemente da qualunque controllore.

\textbf{Nessuno dei tre completa il percorso.} Le metriche fini di questo task
non vanno citate: i tre controllori seguono traiettorie diverse e quindi non
incontrano gli stessi ostacoli. Su questo task il confronto non separa i
controllori, e dirlo \`e il risultato.

\subsection{Impulso laterale: una lacuna comune a tutti e tre}

\begin{table}[htbp]
\centering
\caption{Deviazione laterale residua dopo un disturbo impulsivo (T7).}
\label{tab:t7}
\small
\begin{tabular}{@{}lrrr@{}}
\toprule
 & \Co & \Ct & \Cth \\
\midrule
deviazione residua, impulso $1\times$ [mm]  & 1.34 & 6.04 & 0.164 \\
deviazione residua, impulso $16\times$ [mm] & 393  & 400  & 392 \\
recupera a $16\times$                       & no   & no   & no \\
\bottomrule
\end{tabular}
\end{table}

A impulso grande le tre curve coincidono entro il $2\%$: nessuno dei tre
recupera una frazione misurabile della deviazione. La rilevazione del contatto
non \`e il meccanismo in gioco --- questa \`e una lacuna architetturale comune a
tutti e tre, e la motivazione pi\`u chiara per qualunque controllore futuro su
questa piattaforma.

\subsection{L'inviluppo di stabilit\`a in velocit\`a}
\label{sec:envelope}

Uno sweep di sei fattori di velocit\`a, giudicato sul rimbalzo del corpo con una
soglia dichiarata in anticipo, individua l'inviluppo dell'andatura ad anello
aperto.

\begin{table}[htbp]
\centering
\caption{Rimbalzo del corpo contro fattore di velocit\`a comandato, \Co{}. Soglia
\SI{10}{\milli\meter}, dichiarata prima dello sweep.}
\label{tab:envelope}
\small
\begin{tabular}{@{}lrl@{}}
\toprule
\textbf{fattore} & \textbf{rimbalzo [mm]} & \\
\midrule
1.00 & 3.2  & punto di lavoro nominale \\
1.20 & 3.4  & ultima cella con moto del corpo stabile \\
1.30 & 9.4  & alla soglia \\
1.40 & 19.7 & fallisce \\
\bottomrule
\end{tabular}
\end{table}

\figph{65mm}{fig:envelope}{Sweep di velocit\`a: rimbalzo del corpo e frazione del task contro il fattore di velocit\`a comandato.}{grafico rimbalzo del corpo e frazione del task vs fattore di velocita'}

% =====================================================================
\section{Fattibilit\`a sugli attuatori reali}
\label{sec:feasibility}

La campagna gira ad attuatore ideale: il simulatore concede qualunque coppia.
L'AX-12A stalla a \SI{1.5}{\newton\meter}, e il costruttore raccomanda di
lavorare a un quinto dello stallo per un moto stabile. Poich\'e i giunti sono
comandati in posizione, la coppia \`e una reazione e non un ingresso, quindi
quanto segue misura ci\`o che ciascun controllore \emph{chiede}, non come si
comporterebbe con i motori veri.

\begin{table}[htbp]
\centering
\caption{Coppia comandata contro il limite dell'AX-12A. RMS \`e il giunto pi\`u
caricato di ciascuna run, come frazione dello stallo; ``picco'' \`e la coppia
massima come multiplo dello stallo; ``oltre'' \`e la frazione di campioni sopra
lo stallo.}
\label{tab:feasibility}
\small
\begin{tabular}{@{}lrrrrrrrrr@{}}
\toprule
& \multicolumn{3}{c}{RMS / stallo} & \multicolumn{3}{c}{picco / stallo}
& \multicolumn{3}{c}{campioni oltre} \\
\cmidrule(lr){2-4}\cmidrule(lr){5-7}\cmidrule(lr){8-10}
task & \Co & \Ct & \Cth & \Co & \Ct & \Cth & \Co & \Ct & \Cth \\
\midrule
T2  & 46\% & 44\% & 39\% & 1.2 & 1.3 & 1.5 & 0.40\% & 0.45\% & 0.38\% \\
T3  & 46\% & 42\% & 39\% & 1.3 & 1.4 & 1.6 & 0.50\% & 0.40\% & 0.43\% \\
T4  & 44\% & 44\% & 45\% & 1.4 & 2.3 & 1.8 & 0.39\% & 0.49\% & 0.48\% \\
T4D & 42\% & 42\% & 39\% & 1.7 & 2.7 & 2.1 & 0.45\% & 0.41\% & 0.49\% \\
T5  & 45\% & 44\% & 40\% & 3.0 & 5.6 & 2.7 & 0.61\% & 0.50\% & 0.51\% \\
T6  & 42\% & 43\% & 42\% & 5.0 & 4.9 & 2.8 & 0.91\% & 0.67\% & 0.81\% \\
T7  & 46\% & 44\% & 39\% & 1.2 & 1.3 & 1.5 & 0.41\% & 0.45\% & 0.38\% \\
\bottomrule
\end{tabular}
\end{table}

Due regimi, e vanno letti separatamente:

\begin{itemize}
\item \textbf{La marcia supera il carico raccomandato.} Il giunto pi\`u caricato
sta fra il $39\%$ e il $46\%$ dello stallo su ogni task e per ogni controllore
--- $2.3\times$ il quinto di stallo che il costruttore indica per un moto
stabile. Resta sotto lo stallo, ma non dentro l'inviluppo raccomandato.
\item \textbf{I picchi d'urto escono dallo stallo}, fino a $5.0\times$ sul
percorso a ostacoli, ma su al massimo lo $0.91\%$ dei campioni: sono eventi
isolati, uno per impatto del piede, di durata dell'ordine delle decine di
millisecondi.
\end{itemize}

\Cth{} chiede sistematicamente la coppia continua pi\`u bassa ($39$--$45\%$) e i
picchi d'urto pi\`u bassi sui task con discontinuit\`a, il che \`e coerente con un
anello d'assetto che assorbe parte dell'urto nel moto del corpo invece che nei
giunti.

\figph{65mm}{fig:feasibility}{Coppia comandata contro il limite dell'AX-12A: continua (RMS del giunto peggiore) e di picco, per task e controllore.}{grafico a due pannelli: RMS/limite e picco/limite per task, tre controllori}

% =====================================================================
\section{Discussione}
\label{sec:discussion}

\subsection{I due meccanismi fanno cose diverse, e solo uno paga}

Leggere le Tabelle~\ref{tab:t2t3}--\ref{tab:feasibility} lungo i due confronti
fra adiacenti d\`a un quadro coerente su tutti e sette i task.

\paragraph{La sola ricerca del terreno (\Co$\rightarrow$\Ct) non \`e un
meccanismo d'assetto.} Sui due task in cui il terreno \`e esattamente quello che
il generatore di traiettoria assume --- piano (T2) e curva (T3) --- peggiora il
beccheggio di tre o quattro volte, con un margine da $16$ a $40$ volte il
rumore di fondo, e non restituisce nulla di misurabile in energia. Sui task con
terreno reale fa meglio: il rollio migliora del $30\%$ sul dosso e del $49\%$
sull'ostacolo, il costo di trasporto cala del $7$--$10\%$ e la velocit\`a media
sale del $2$--$3\%$. Ma il beccheggio, che \`e il grado di libert\`a su cui la
ricerca agisce, peggiora comunque del $21\%$ sul dosso.

La ragione \`e visibile nel meccanismo pi\`u che nei numeri. La ricerca abbassa
un piede in volo finch\'e il residuo di coppia non dichiara il contatto; la
dichiarazione \`e fatta per zampa, in modo asincrono, e la correzione di altezza
che ne risulta \`e un gradino. Su un terreno che il generatore gi\`a modella
correttamente i gradini sono spuri, e tre zampe che gradinano a istanti
leggermente diversi sono esattamente un'eccitazione di beccheggio. Arrigoni et
al.~\cite{arrigoni2024} introducono il meccanismo per terreno che il robot non
pu\`o vedere; la nostra misura dice che su terreno che \emph{pu\`o} vedere il
meccanismo \`e un costo netto, e che quel costo non scompare quando il terreno
diventa irregolare --- semplicemente gli si affianca un beneficio sugli altri
assi.

\paragraph{La retroazione d'assetto (\Ct$\rightarrow$\Cth) \`e il meccanismo che
compra assetto, e lo compra con l'energia.} Il segno \`e lo stesso su ogni task
e ogni ampiezza supera largamente la regola del $3\times$: beccheggio $-82\%$
(piano), $-79\%$ (curva), $-81\%$ (dosso), $-73\%$ (ostacolo), con rapporti sul
rumore da $18$ a $94$. Sulla rampa di $8^\circ$ la differenza \`e qualitativa
anzich\'e incrementale: \Co{} e \Ct{} portano il tronco parallelo alla rampa
($0.11^\circ$ e $0.26^\circ$ di angolo corpo-meno-rampa), mentre \Cth{} lo
tiene entro mezzo grado dall'\emph{orizzontale} ($-7.54^\circ$ di
corpo-meno-rampa su una pendenza di $8^\circ$). Sono due comportamenti diversi,
non due qualit\`a dello stesso comportamento, ed \`e per questo che quella riga
non porta verdetti percentuali.

Il prezzo \`e altrettanto costante: $+28\%$ di energia e costo di trasporto su
piano e curva, $+40\%$ su dosso e rampa, $+30\%$ sull'ostacolo, tutti da $17$ a
$28$ volte il rumore di fondo. \`E interessante che la coppia \emph{continua}
scenda ($-3$ fino a $-8\%$, dichiarabile su cinque task): l'energia in pi\`u \`e
spesa nella corsa verticale aggiuntiva dei piedi, non nel sostenere carichi di
giunto maggiori.

\paragraph{La composizione non \`e la somma.} \Cth{} \`e \Ct{} con l'anello
d'assetto in serie dopo la ricerca del terreno, quindi il $+279\%$ di
beccheggio della ricerca e il $-82\%$ dell'anello sono misurati sullo stesso
percorso di segnale. Confrontando i due estremi, \Cth{} finisce \emph{meglio di
\Co} sul beccheggio ($0.356$ contro $0.514$ deg in piano), cio\`e l'anello
d'assetto ripaga pi\`u che interamente il disturbo introdotto dalla ricerca. Se
un controllore fatto di sola retroazione d'assetto \emph{senza} ricerca del
terreno farebbe ancora meglio \`e l'ablazione successiva ovvia, e non l'abbiamo
eseguita: l'impianto per farlo non esiste, e costruirlo avrebbe significato
modificare un modello gi\`a congelato per la campagna.

\subsection{Un risultato sulla composizione stessa: la rilevazione
dell'appoggio si degrada}

Il risultato architetturale pi\`u chiaro del progetto non era pianificato.
Quando il task della rampa (T4) \`e stato analizzato su \Cth, il rilevatore
cinematico di appoggio (Sezione~\ref{sec:setup}) --- un piede \`e in appoggio
quando la sua velocit\`a nel mondo resta sotto \SI{0.05}{\meter\per\second} per
almeno un quarto dell'appoggio nominale --- non trovava praticamente alcun
appoggio. La nostra prima lettura \`e stata che il robot non camminasse. Era
sbagliata, e vale la pena conservare il perch\'e: la quota media del corpo che
avevamo citato era l'altezza della rampa stessa, e il conteggio di appoggi dai
sensori di forza pari a $1.00$ non ha significato su una rampa, perch\'e quei
sensori vedono solo il piano del pavimento.

La lettura corretta \`e venuta dal contare, per zampa, la frazione di campioni
sotto soglia e il tratto consecutivo pi\`u lungo:

\begin{center}
\small
\begin{tabular}{@{}lrr@{}}
\toprule
 & campioni sotto \SI{0.05}{\meter\per\second} & tratto consecutivo pi\`u lungo \\
\midrule
\Ct  & 41--54\% & 71--113 campioni ($\approx$ un appoggio intero) \\
\Cth & 4--8\%   & 3--12 campioni \\
\bottomrule
\end{tabular}
\end{center}

I piedi non mancano l'atterraggio: atterrano e poi vengono rimossi. L'anello
d'assetto aggiunge una correzione continua di altezza
$\Delta z = x\,u_\theta + y\,u_\phi$ a ogni piede, compresi quelli in appoggio,
quindi un piede in appoggio non \`e mai fermo nel riferimento mondo.
L'ampiezza \`e piccola --- con una velocit\`a mediana del piede di
\SI{0.14}{\meter\per\second} il moto su un appoggio frammentato \`e circa
\SI{17}{\milli\meter}, contenuto dalla saturazione e dal limitatore di
pendenza --- quindi l'andatura non ne risente. Ci\`o che si rompe \`e
l'\emph{osservabilit\`a del contatto da parte di qualunque criterio basato sulla
quiete}.

Questo conta al di l\`a del nostro banco, perch\'e la ricerca del terreno
sottostante \`e essa stessa un meccanismo di rilevazione del contatto. Comporre
un anello d'assetto in serie dopo un meccanismo innescato dal contatto d\`a al
secondo un mondo in cui nulla \`e mai fermo. La composizione funziona qui perch\'e
la ricerca del terreno usa un residuo di coppia e non la cinematica del piede,
e perch\'e la correzione d'assetto \`e limitata in pendenza; non sopravvivrebbe a
un rilevatore di appoggio costruito sulla velocit\`a del piede. Consideriamo
questa la cosa pi\`u trasferibile prodotta dal progetto.

\subsection{Che cosa la batteria di task non ha potuto risolvere}

Tre dei sette task non hanno dato verdetto, e dirlo fa parte del risultato.

\begin{itemize}
\item \textbf{T6 (percorso a sette ostacoli) non discrimina.} Nessuno dei tre
controllori lo completa: $2.91$, $3.12$, \SI{2.99}{\meter} percorsi, tutti e
tre sotto la soglia di validit\`a sulla deriva laterale (\SI{0.183}{\meter},
ricavata dalla geometria del percorso e non dai controllori). Oltre quella
deriva il robot pu\`o essere passato \emph{a fianco} dell'ostacolo~3 invece che
sopra, quindi n\'e la distanza n\'e alcuna metrica a valle sono confrontabili: i
tre controllori non hanno incontrato gli stessi ostacoli. Questo \`e un
cambiamento rispetto a una campagna precedente, prima che il limitatore di
pendenza fosse presente su tutti e tre gli impianti, in cui \Ct{} sembrava
arrivare a \SI{3.71}{\meter} contro \SI{2.44}{\meter} di \Co. Quel titolo \`e
superato e non va citato.

\item \textbf{T7 (impulso laterale) separa solo alla piccola ampiezza.} A
$1\times$ la deviazione residua \`e $1.34$, $6.04$, \SI{0.164}{\milli\meter} ---
un miglioramento del $97\%$ per \Cth{} --- ma il rumore di fondo non \`e stato
misurato su questo task, quindi non si dichiara alcun verdetto. A $16\times$
tutti e tre finiscono a $392$--\SI{400}{\milli\meter} e nessuno recupera: il
disturbo \`e semplicemente pi\`u grande di quanto ciascuno dei tre possa
reiettare.

\item \textbf{La coppia di picco non supera mai la regola del $3\times$}, su
nessun task e nessun confronto, nonostante differenze fino al $+87\%$. Il picco
\`e una statistica a campione singolo di un urto e la sua dispersione fra run \`e
dello stesso ordine delle differenze fra controllori. La colonna RMS, che \`e un
integrale, \`e dichiarabile su cinque task. Riportiamo entrambe e traiamo
conclusioni solo dalla seconda.
\end{itemize}

% =====================================================================
\section{Limiti}
\label{sec:limits}

\begin{description}[style=unboxed,leftmargin=0pt,font=\normalfont\itshape]

\item[Attuatori ideali.] I giunti sono comandati in posizione e il simulatore
fornisce qualunque coppia il comando richieda. La
Sezione~\ref{sec:feasibility} misura quindi la \emph{domanda}, non il
comportamento ottenibile. Una campagna con i limiti di velocit\`a e coppia
dell'AX-12A nell'anello comprimerebbe probabilmente le differenze, perch\'e
\Cth{} \`e il controllore che chiede la maggiore corsa aggiuntiva del piede per
unit\`a di tempo.

\item[I sensori di contatto vedono solo il pavimento.] I sensori di forza
Simscape sono attaccati al piano del terreno, quindi su rampa, dosso e ostacoli
non riportano nulla. Ogni metrica dipendente dall'appoggio su T4, T4D, T5 e T6
viene invece dal rilevatore cinematico --- che, come mostra la
Sezione~\ref{sec:discussion}, \`e a sua volta degradato su \Cth. Dove i due
rilevatori sono in disaccordo riportiamo il disaccordo invece di sceglierne
uno.

\item[I bracci di leva dell'assetto sono cablati.] Le $x$ e $y$ usate in
$\Delta z = x\,u_\theta + y\,u_\phi$ sono costanti nel modello, circa
$1.8\times$ pi\`u piccole delle posizioni nominali dei piedi
\texttt{cfg.pf\_nom}, pur con segni e ordinamento delle zampe corretti. Abbiamo
scelto di \emph{non} correggerle: i guadagni proporzionali sono stati tarati
empiricamente con quei valori nell'anello, quindi cambiare i bracci di leva
senza ritarare cambierebbe il guadagno d'anello dello stesso fattore. La coppia
(braccio di leva, guadagno) \`e coerente; solo la sua fattorizzazione \`e
arbitraria.

\item[Lo stimatore di dinamica inversa non \`e l'impianto.] Dopo la correzione
della Sezione~\ref{sec:estimator} il disallineamento residuo \`e l'\SI{1.6}{\percent}
di massa nei piedi, che l'albero URDF non porta. La soglia di contatto
$\sigma_j$ lo assorbe, ma \`e stata tarata su questo banco e andrebbe ritarata
su un altro.

\item[La soglia \`e autotarata.] $\sigma_j$ \`e stata fissata a partire dal
residuo osservato durante un volo notoriamente libero su questo impianto. Non
\`e quindi un parametro indipendente, e alla ricerca del terreno non si pu\`o
accreditare robustezza rispetto a una soglia su cui \`e stata adattata.

\item[Solo stabilit\`a quasi-statica.] Il numero di Froude della marcia nominale
\`e $v^2/(gh) = 0.0095$, quindi lo ZMP collassa sulla proiezione a terra del
centro di massa entro la precisione della misura; l'analisi di stabilit\`a
assume una zona di contatto planare e riporta l'inclinazione del piano
d'appoggio, in modo che l'ipotesi possa essere verificata invece che data per
buona.

\item[Nessun hardware.] Nulla in questo report \`e stato misurato su un PhantomX
fisico. La piattaforma, l'URDF e il datasheet del servo sono reali; il
confronto non \`e validato.

\end{description}

% =====================================================================
\section{Contributi per componente del gruppo}
\label{sec:contributions}

\noindent\fcolorbox{red!50}{red!3}{\begin{minipage}{0.97\linewidth}
\small\textbf{DA FARE --- da confermare e completare a cura degli autori prima
della consegna.} La suddivisione qui sotto riflette la divisione del lavoro
come \`e effettivamente avvenuta; nomi e qualunque voce che uno dei due voglia
rivendicare vanno sistemati a mano.
\end{minipage}}

\vspace{6pt}

\begin{description}[style=unboxed,leftmargin=0pt,font=\normalfont\bfseries]

\item[Studente 1 --- impianto e controllori.]
Costruzione dell'impianto Simscape Multibody a partire dall'URDF del PhantomX;
implementazione del generatore a tripode ad anello aperto (\Co);
implementazione dello strato di ricerca del terreno secondo Arrigoni et al.\
(\Ct), compresi il residuo di dinamica inversa e la flag di contatto per zampa;
progetto, implementazione e taratura empirica dell'anello d'assetto (\Cth) ---
prima il beccheggio, poi il rollio --- e della saturazione e del limitatore di
pendenza che ne vincolano l'autorit\`a. Proprietà dei modelli Simulink per tutto
il progetto.

\item[Studente 2 --- misura, diagnosi e analisi.]
Progetto della batteria di sette task e del protocollo di misura
(Sezione~\ref{sec:method}), compresi la procedura per il rumore di fondo e la
regola di dichiarabilit\`a a $3\times$; l'intera campagna di misura e la sua
automazione (scelta del controllore in un punto unico, esecuzione batch,
conversione dei risultati, generazione della tabella di confronto);
individuazione, quantificazione e correzione dei due difetti del banco della
Sezione~\ref{sec:defects}; la diagnosi sulla rilevazione dell'appoggio della
Sezione~\ref{sec:discussion}; la verifica di stabilit\`a quasi-statica; l'analisi
di fattibilit\`a rispetto al datasheet del servo; struttura e documentazione del
repository.

\item[In comune.]
Scelta dei tre controllori e della struttura dell'ablazione; ricostruzione e
giustificazione dei parametri d'andatura (Sezione~\ref{sec:setup});
interpretazione dei risultati; questo report.

\end{description}

% =====================================================================
\section{Rapporto con altri corsi e novit\`a}
\label{sec:novelty}

Questo progetto non deriva da, e non estende, un progetto presentato per un
altro corso. Il modello della piattaforma, i tre controllori, la batteria di
task, il protocollo di misura e tutti i risultati qui riportati sono stati
prodotti per Field and Service Robotics.

Rispetto al materiale del corso il progetto \`e un'applicazione pi\`u che
un'estensione della teoria, e la sua novit\`a sta in tre punti:

\begin{enumerate}
\item \textbf{Un'ablazione controllata invece di una dimostrazione.} La forma
abituale di un progetto di questo tipo \`e: implementare un controllore,
mostrare che cammina. Qui lo stesso impianto, terreno, legge di contatto e
generatore d'andatura portano tre controllori che differiscono per esattamente
un meccanismo ciascuno, quindi ogni confronto isola un meccanismo. \`E quel
disegno a rendere possibile affermare che la ricerca del terreno costa assetto
e che la retroazione d'assetto costa energia, invece che affermare che un
controllore completo \`e migliore di un altro.

\item \textbf{Un criterio di dichiarabilit\`a fissato prima della campagna.} I
risultati di simulazione sono riproducibili all'ultima cifra se non si perturba
nulla, il che fa sembrare significativa ogni differenza. Perturbare l'altezza
iniziale di \SI{\pm 0.2}{\milli\meter} su tre run d\`a un rumore di fondo, e
richiedere che una differenza superi $3\times$ il peggiore dei due rumori prima
di dichiararla converte una tabella di numeri in una tabella di affermazioni.
Circa un terzo delle differenze misurate non sopravvive a quella regola,
compresi tutti i confronti sulla coppia di picco.

\item \textbf{Due difetti silenziosi del banco, trovati per misura e non per
ispezione.} Entrambi (Sezione~\ref{sec:defects}) producevano numeri plausibili
e dall'aria pubblicabile pur essendo sbagliati, e nessuno dei due era visibile
nello schema a blocchi. Quello del disallineamento dello stimatore in
particolare \`e un'istanza diretta di ci\`o che il corso dice sui residui basati
su modello: correggere l'impianto senza correggere il modello che lo prevede \`e
peggio che non correggere nessuno dei due, perch\'e il residuo satura e la flag
di contatto degenera in una costante.
\end{enumerate}

% =====================================================================
\section{Conclusioni}
\label{sec:conclusions}

Tre controllori per un esapode PhantomX sono stati confrontati su un unico
impianto lungo sette task, con un protocollo che dichiara una differenza solo
quando supera tre volte il rumore di fondo misurato.

\begin{itemize}
\item La ricerca del terreno di~\cite{arrigoni2024}, presa da sola, \`e un costo
netto su terreno che il generatore d'andatura gi\`a modella: il beccheggio
peggiora del $219$--$279\%$ senza alcun ritorno energetico. Su terreno
irregolare migliora il rollio ($-30$ fino a $-49\%$), il costo di trasporto
($-7$ fino a $-10\%$) e la velocit\`a ($+2$ fino a $+3\%$), pur continuando a
peggiorare il beccheggio.
\item L'anello d'assetto \`e il meccanismo che compra assetto, e lo fa su ogni
task: beccheggio da $-73$ a $-82\%$ a un rapporto da $18$ a $94$ volte il
rumore di fondo, e sulla rampa di $8^\circ$ un comportamento qualitativamente
diverso (tronco tenuto orizzontale invece che parallelo alla pendenza). Costa
il $28$--$40\%$ in energia e costo di trasporto, speso in corsa dei piedi e non
in carico ai giunti --- la coppia continua anzi cala.
\item Composti in serie, i due meccanismi lasciano il corpo migliore del
riferimento ad anello aperto sul beccheggio, ma rendono il contatto del piede
inosservabile a qualunque rilevatore basato sulla quiete: i tratti di appoggio
collassano da $71$--$113$ campioni a $3$--$12$.
\item Sui servo reali la marcia da sola girerebbe al $39$--$46\%$ dello stallo
sul giunto pi\`u caricato --- $2.3\times$ il carico continuo raccomandato dal
costruttore --- con picchi d'urto sopra lo stallo su al massimo lo $0.91\%$ dei
campioni.
\end{itemize}

La continuazione ovvia \`e la quarta cella dell'ablazione --- retroazione
d'assetto senza ricerca del terreno --- che direbbe se la ricerca contribuisca
qualcosa una volta che il corpo \`e stabilizzato, e una campagna con i limiti
dei servo nell'anello, che \`e dove il costo energetico di \Cth{} smetterebbe di
essere gratuito.

% =====================================================================
\begin{thebibliography}{9}

\bibitem{arrigoni2024}
S.~Arrigoni, F.~Braghin, et al.,
``Control of a Hexapod Robot Considering Terrain Interaction,''
\emph{Robotics}, vol.~13, no.~10, art.~142, 2024.
% TODO: verificare la lista completa degli autori sul PDF

\bibitem{ax12a}
ROBOTIS, \emph{Dynamixel AX-12A e-Manual}, ROBOTIS Co., Ltd.
\url{https://emanual.robotis.com/docs/en/dxl/ax/ax-12a/}

\bibitem{phantomx}
Trossen Robotics, \emph{PhantomX AX Metal Hexapod Mark~II --- assembly and
specifications}.

\bibitem{simscape}
The MathWorks, Inc., \emph{Simscape Multibody User's Guide}, R2025b.

\end{thebibliography}

\end{document}
